
\documentclass{article}
\usepackage{colortbl}
\usepackage{makecell}
\usepackage{multirow}
\usepackage{supertabular}

\begin{document}

\newcounter{utterance}

\centering \large Interaction Transcript for game `cladder', experiment `full\_v1.5\_default', episode 8727 with refineonly{-}Qwen3.5{-}9B.
\vspace{24pt}

{ \footnotesize  \setcounter{utterance}{1}
\setlength{\tabcolsep}{0pt}
\begin{supertabular}{c@{$\;$}|p{.15\linewidth}@{}p{.15\linewidth}p{.15\linewidth}p{.15\linewidth}p{.15\linewidth}p{.15\linewidth}}
    \# & \multicolumn{2}{c}{Player} && \multicolumn{2}{c}{Game Master} \\
    \hline

    \theutterance \stepcounter{utterance}  
    & & & \multicolumn{4}{p{0.6\linewidth}}{
        \cellcolor[rgb]{0.9,0.9,0.9}{
            \makecell[{{p{\linewidth}}}]{
                \texttt{\tiny{[P1$\langle$GM]}}
                \texttt{You are an expert in causal inference. The following question is not a typical commonsense query, but rather a meticulously designed question created by a professor specializing in causal inference, intended to assess the students' mastery of the course content.} \\
\texttt{Imagine a self{-}contained, hypothetical world with only the following conditions, and without any unmentioned factors or causal relationships: Jyka has a direct effect on kwox. Yupt has a direct effect on kwox. For those who are not jyka, the probability of kwox is 32\%. For those who are jyka, the probability of kwox is 38\%. Will jyka increase the chance of kwox?} \\
\texttt{Start your answer with 'Yes' or 'No', followed by additional reasoning or evidence to support your explanation.} \\
            }
        }
    }
     \\ \\

    \theutterance \stepcounter{utterance}  
    & & \multicolumn{4}{p{0.6\linewidth}}{
        \cellcolor[rgb]{0.95,0.95,0.95}{
            \makecell[{{p{\linewidth}}}]{
                \texttt{\tiny{[P1$\leftrightarrow$P]}}
                \texttt{{[}DRAFT{]}} \\
\texttt{Yes} \\
\\ 
\texttt{**Reasoning:**} \\
\\ 
\texttt{In this hypothetical world, we are evaluating the causal effect of **Jyka** on **kwox**. The problem explicitly states that Jyka has a *direct effect* on kwox and provides the conditional probabilities for the outcome \$P(\textbackslash{}text\{kwox\} | \textbackslash{}text\{Jyka\})\$ and \$P(\textbackslash{}text\{kwox\} | \textbackslash{}neg\textbackslash{}text\{Jyka\})\$.} \\
\\ 
\texttt{1.  **Identify the Counterfactuals**:} \\
\texttt{    *   The probability of kwox occurring when Jyka is present (the "treated" state) is given as \$P(\textbackslash{}text\{kwox\} | \textbackslash{}text\{Jyka\}) = 38\textbackslash{}\%\$.} \\
\texttt{    *   The probability of kwox occurring when Jyka is absent (the "control" state) is given as \$P(\textbackslash{}text\{kwox\} | \textbackslash{}neg\textbackslash{}text\{Jyka\}) = 32\textbackslash{}\%\$.} \\
\\ 
\texttt{2.  **Calculate the Average Causal Effect (ACE)**:} \\
\texttt{    Since the prompt specifies a self{-}contained world with no unmentioned factors or confounders, we can directly compare these conditional probabilities to determine the net effect of Jyka. The causal effect is the difference between the outcome probability under intervention (Jyka) and the outcome probability under non{-}intervention (no Jyka).} \\
\texttt{    \$\$ \textbackslash{}text\{Effect\} = P(\textbackslash{}text\{kwox\} | \textbackslash{}text\{Jyka\}) {-} P(\textbackslash{}text\{kwox\} | \textbackslash{}neg\textbackslash{}text\{Jyka\}) \$\$} \\
\texttt{    \$\$ \textbackslash{}text\{Effect\} = 38\textbackslash{}\% {-} 32\textbackslash{}\% = +6\textbackslash{}\% \$\$} \\
\\ 
\texttt{3.  **Conclusion**:} \\
\texttt{    Because the probability of kwox increases from 32\% to 38\% when Jyka is introduced, Jyka causally increases the chance of kwox by 6 percentage points. The presence of Yupt as another direct cause does not alter this specific comparison regarding Jyka's contribution, as the probabilities provided already account for the state of the world in each condition.} \\
\\ 
\texttt{Therefore, Jyka increases the chance of kwox.} \\
            }
        }
    }
    & \\ \\

    \theutterance \stepcounter{utterance}  
    & \multicolumn{4}{p{0.6\linewidth}}{
        \cellcolor[rgb]{0.561,0.737,0.561}{
            \makecell[{{p{\linewidth}}}]{
                \texttt{\tiny{[P1$\rangle$GM]}}
                \texttt{Yes} \\
\\ 
\texttt{**Reasoning:**} \\
\\ 
\texttt{In this hypothetical world, we are evaluating the causal effect of **Jyka** on **kwox**. The problem explicitly states that Jyka has a *direct effect* on kwox and provides the conditional probabilities for the outcome \$P(\textbackslash{}text\{kwox\} | \textbackslash{}text\{Jyka\})\$ and \$P(\textbackslash{}text\{kwox\} | \textbackslash{}neg\textbackslash{}text\{Jyka\})\$.} \\
\\ 
\texttt{1.  **Identify the Counterfactuals**:} \\
\texttt{    *   The probability of kwox occurring when Jyka is present (the "treated" state) is given as \$P(\textbackslash{}text\{kwox\} | \textbackslash{}text\{Jyka\}) = 38\textbackslash{}\%\$.} \\
\texttt{    *   The probability of kwox occurring when Jyka is absent (the "control" state) is given as \$P(\textbackslash{}text\{kwox\} | \textbackslash{}neg\textbackslash{}text\{Jyka\}) = 32\textbackslash{}\%\$.} \\
\\ 
\texttt{2.  **Calculate the Average Causal Effect (ACE)**:} \\
\texttt{    Since the prompt specifies a self{-}contained world with no unmentioned factors or confounders, we can directly compare these conditional probabilities to determine the net effect of Jyka. The causal effect is the difference between the outcome probability under intervention (Jyka) and the outcome probability under non{-}intervention (no Jyka).} \\
\texttt{    \$\$ \textbackslash{}text\{Effect\} = P(\textbackslash{}text\{kwox\} | \textbackslash{}text\{Jyka\}) {-} P(\textbackslash{}text\{kwox\} | \textbackslash{}neg\textbackslash{}text\{Jyka\}) \$\$} \\
\texttt{    \$\$ \textbackslash{}text\{Effect\} = 38\textbackslash{}\% {-} 32\textbackslash{}\% = +6\textbackslash{}\% \$\$} \\
\\ 
\texttt{3.  **Conclusion**:} \\
\texttt{    Because the probability of kwox increases from 32\% to 38\% when Jyka is introduced, Jyka causally increases the chance of kwox by 6 percentage points. The presence of Yupt as another direct cause does not alter this specific comparison regarding Jyka's contribution, as the probabilities provided already account for the state of the world in each condition.} \\
\\ 
\texttt{Therefore, Jyka increases the chance of kwox.} \\
            }
        }
    }
    & & \\ \\

    \theutterance \stepcounter{utterance}  
    & & & \multicolumn{2}{p{0.3\linewidth}}{
        \cellcolor[rgb]{0.9,0.9,0.9}{
            \makecell[{{p{\linewidth}}}]{
                \texttt{\tiny{[GM$|$GM]}}
                \texttt{Yes} \\
            }
        }
    }
    & & \\ \\

    \theutterance \stepcounter{utterance}  
    & & & \multicolumn{2}{p{0.3\linewidth}}{
        \cellcolor[rgb]{0.9,0.9,0.9}{
            \makecell[{{p{\linewidth}}}]{
                \texttt{\tiny{[GM$|$GM]}}
                \texttt{yes} \\
            }
        }
    }
    & & \\ \\

    \theutterance \stepcounter{utterance}  
    & & & \multicolumn{2}{p{0.3\linewidth}}{
        \cellcolor[rgb]{0.9,0.9,0.9}{
            \makecell[{{p{\linewidth}}}]{
                \texttt{\tiny{[GM$|$GM]}}
                \texttt{game\_result = WIN} \\
            }
        }
    }
    & & \\ \\

\end{supertabular}
}

\end{document}
